% %%%%%%%%%%%%%%%%%%%%%%%%%%%%%%%%%%%%%%%%%%%%%%%%%%%%%%%%%%%%%%%%%%%%%% %
%                                                                         %
% The Project Gutenberg EBook of a Test Book, by A. Fixture               %
%                                                                         %
% Title: A Test Book Shaped Like Calculus Made Easy                       %
%        Being a fixture for EDIE                                         %
%                                                                         %
% Author: A. Fixture                                                      %
%                                                                         %
% Release Date: November 14, 2026 [EBook #99999]                          %
%                                                                         %
% Language: English                                                       %
%                                                                         %
% Character set encoding: UTF-8                                           %
% %%%%%%%%%%%%%%%%%%%%%%%%%%%%%%%%%%%%%%%%%%%%%%%%%%%%%%%%%%%%%%%%%%%%%% %
\documentclass[12pt]{book}
\usepackage{amsmath}
\usepackage{graphicx}
\usepackage{fancyvrb}
\DefineVerbatimEnvironment{PGtext}{Verbatim}{fontsize=\small}

\makeatletter
% \Chapter[running head]{number}{title}; the running head is optional.
\newcommand{\Chapter}{\@ifnextchar[{\@Chapter}{\@@Chapter}}
\def\@Chapter[#1]#2#3{\chapter*{#2. #3}\markboth{#1}{#1}}
\def\@@Chapter#1#2{\@Chapter[#2]{#1}{#2}}
\makeatother
\newcommand{\ChapterStar}[2][]{\chapter*{#2}\markboth{#1}{#1}}
\newcommand{\AltChapter}[3][]{\chapter*{#2. #3}}
\newcommand{\Section}[1]{\section*{#1}}
\newcommand{\Subsection}[1]{\subsection*{#1}}
\newcommand{\Paragraph}[1]{\paragraph*{#1}}
\newcommand{\Case}[1]{\par\textit{Case #1.}}
\newcommand{\Example}[1]{\par\textit{Example #1.}}
\newcommand{\Figure}[2][]{\begin{center}\includegraphics[#1]{./images/#2}\end{center}}
\newcommand{\Graphic}[2]{\includegraphics[width=#1]{./images/#2}}
\newcommand{\Pagelabel}[1]{\label{page:#1}}

\title{A Test Book}
\author{A. Fixture}

\begin{document}

\begin{PGtext}
The Project Gutenberg EBook. Braces here are not LaTeX: { and } and \Chapter{IX}{Not a Chapter}
\end{PGtext}

\ChapterStar[Prologue]{Prologue}

What one fool can do, another can.  Costs 5\% more, with \{braces\} and \$1.

\Chapter{I}{To Deliver You from the Preliminary Terrors.}

The preliminary terror is the symbol $dx$, and the sign $\int$ which is merely
a long~$S$.
% \Chapter{X}{A Chapter That Is Only a Comment} with a stray } brace
Thus $\int dx$ means the sum of all the little bits of~$x$.

\Chapter[On Different Degrees of Smallness]{II}{On Different Degrees of\\
Smallness.}

\Section{Small quantities of the {\itshape very} {first {order}}}
We shall find that
\[
  \frac{dy}{dx} = \frac{\alpha + \beta}{\gamma}
\]
continues the same paragraph, and $$\int_0^1 x\,dx = \tfrac{1}{2}$$ does too.

\Chapter[Relative Growing]
{III}{On Relative Growings.}

\Subsection{Growing quantities}
\Paragraph{Case of a square.}
Let $x$ grow.
\Case{1}
Here
\begin{align*}
  y + dy &= (x + dx)^2 \\
  dy &= 2x\,dx + (dx)^2.
\end{align*}

\Example{(1)}
Differentiate $y = x^{n}$.
\Figure[2in]{fig1}

A small diagram \Graphic{0.5in}{fig2} sits in the text, and one sits in mathematics:
$\Graphic{0.3in}{fig3} = \dfrac{dy}{dx}$.

\begin{figure}[htbp]
  \centering
  \includegraphics[width=2in]{images/fig4}
  \caption{Fig.~4. A curve and its slope.}
  \label{fig:curve}
\end{figure}

\Figure{fig99}

\AltChapter[Alt]{XV}{How to Deal with Sines and Cosines}
\Pagelabel{162}
\begin{equation}
  \frac{d(\sin\theta)}{d\theta} = \cos\theta
\end{equation}

\end{document}
